% SEGMENT OLP-0614-S01
% Part: methods
% Chapter: induction
% Section: introduction

\documentclass[../../../include/open-logic-section]{subfiles}

\begin{document}

\olfileid{mth}{ind}{int}

\olsection{அறிமுகம்}

தருக்கவியல், கோட்பாட்டுக் கணினியியல், கணிதம் ஆகியவற்றின்
ஏறக்குறைய எல்லாத் துறைகளிலும் வெவ்வேறு வடிவங்களில்
பயன்படும் முக்கியமான நிறுவல் உத்தி தொகுத்தறிதல்.
தருக்கவியலின் பல முடிவுகளை நிறுவ இது தேவைப்படுகிறது.

% SEGMENT OLP-0614-S02
தொகுத்தறிதல் பெரும்பாலும் வருவித்தலுக்கு எதிராக
வைக்கப்பட்டு, குறிப்பிட்ட நிகழ்வுகளிலிருந்து பொதுக்
கூற்றுக்கு இட்டுச் செல்லும் உய்த்தறிதல் என்று
விளக்கப்படுகிறது. எடுத்துக்காட்டாக, பல மரகதக் கற்கள்
பச்சையாக இருப்பதையும், பச்சையாக இல்லாத மரகதக் கல்
எதையும் காணாததையும் வைத்து, எல்லா மரகதக் கற்களும்
பச்சை என்று முடிவு செய்யலாம். இது தொகுத்தறிதல் வழியான
உய்த்தறிதல்: இந்த மரகதக் கல் பச்சை, அந்த மரகதக் கல்
பச்சை போன்ற பல குறிப்பிட்ட நிகழ்வுகளிலிருந்து,
எல்லா மரகதக் கற்களும் பச்சை என்ற பொதுக் கூற்றுக்கு
இட்டுச் செல்கிறது. \emph{கணிதத்} தொகுத்தறிதலும்
ஒரு பொதுக் கூற்றை முடிவாகக் கொள்ளும் உய்த்தறிதல்தான்;
ஆனால் இந்த ``எளிய தொகுத்தறிதலை'' விட மிகவும்
வேறுபட்ட வகையைச் சேர்ந்தது.

% SEGMENT OLP-0614-S03
தோராயமாகச் சொன்னால், கணிதத்தில் தொகுத்தறிதல் நிறுவல்,
ஒரு வகையைச் சேர்ந்த எல்லாக் கணிதப் பொருள்களுக்கும்
ஒரு குறிப்பிட்ட பண்பு உண்டு என்று முடிவு செய்கிறது.
மிக எளிய நிலையில், அது இயல் எண்களைப் பற்றியது.
அப்போது எல்லா இயல் எண்களுக்கும் அந்தப் பண்பு உண்டு
என்று நிறுவப் பயன்படுகிறது. அதற்காகக் காட்டுவது:
\begin{enumerate}
    \item $0$ க்கு அந்தப் பண்பு உண்டு; மேலும்
    \item ஓர் எண்~$k$ க்கு அந்தப் பண்பு இருக்கும்
    போதெல்லாம், $k+1$ க்கும் அது உண்டு.
\end{enumerate}

% SEGMENT OLP-0614-S04
இயல் எண்கள் மீதான தொகுத்தறிதலை, எண்களை ஒதுக்கக்கூடிய
பிற கணிதப் பொருள்கள் பற்றிய பொதுக் கூற்றுகளை நிறுவவும்
பல நேரங்களில் பயன்படுத்தலாம். எடுத்துக்காட்டாக,
ஒவ்வொரு முடிவுறு கணத்திற்கும் அதன் உறுப்புகளின்
எண்ணிக்கையாக ஒரு முடிவுறு எண்~$n$ உள்ளது.
ஒவ்வோர் எண்~$n$ க்கும் ``அளவு~$n$ உடைய
எல்லா முடிவுறு கணங்களும் \dots'' என்ற பண்பு உண்டு
என்று தொகுத்தறிதலால் காட்ட முடிந்தால், எல்லா
முடிவுறு கணங்களையும் பற்றிய ஒன்றைக் காட்டியிருப்போம்.

% SEGMENT OLP-0614-S05
\emph{தொகுத்தறிதல் முறையில் வரையறுக்கப்பட்ட} கணிதப்
பொருள்களுக்கும் தொகுத்தறிதலைப் பொதுமைப்படுத்தலாம்.
எடுத்துக்காட்டாக, முதல்தரத் தருக்கத்தில் வருபவை போன்ற
முறைமொழிக் கோவைகள் தொகுத்தறிதல் முறையில்
வரையறுக்கப்படுகின்றன. அத்தகைய எல்லாக் கோவைகளையும்
பற்றிய முடிவுகளை நிறுவ \emph{கட்டமைப்புத் தொகுத்தறிதல்}
ஒரு வழி. குறிப்பாக, அது தருக்கவியலில் மிகவும்
பயனுள்ளது; பரவலாகப் பயன்படுத்தப்படுகிறது.

\end{document}
